\cotefile{74}
\scanpage{73}

% top, lines 1–2
\[
\begin{aligned}
\unsure{120|720}\quad\mathrm{Hex}&=\text{ens des hexagones de trièdres \inserted{liées (de $\Delta$)} ( ``grands hexagones'')}\\
&\simeq\text{ens des pavés}
\end{aligned}
\]

% top, lines 3–4
\[
\begin{aligned}
720\quad\mathrm{hex}&=\text{ens des hexagones de sommets \inserted{de $\Delta$} \deleted{\unsure{$\Delta$}} (= ``petits hex.'')}\\
21\quad &\simeq\text{ens des paires de trièdres opposables}
\end{aligned}
\]

% top, lines 5–6
\[
\begin{aligned}
&\deleted{\simeq\text{ens des couples }(X,t),\ X\text{ pavé et }t\in X}\\
720\quad\mathrm{tria}&=\deleted{\text{ens des \unsure{trièdres|triades}}}
\end{aligned}
\]

% upper, lines 7–8
\[
40\quad\mathrm{triHex}\simeq\text{ens des triples (non ordonnés) \inserted{``distingués''} de grands hexagones,\,/ i.e.}
\]

% upper, lines 9–10
\[
\begin{aligned}
360\quad(\mathrm{tria}')&\simeq\text{ens des \inserted{paires de} trièdres \inserted{opposés} \deleted{\unsure{(liés) = ens des pai|(Tria' = ens des pai}}}\\
&\hspace{2em}\text{(ce sont les diagonales des \uline{grands} Hexagones}
\end{aligned}
\]

% upper, lines 11–12
\[
\begin{aligned}
240\quad\deleted{\unsure{\mathrm{triHex}}}\inserted{\mathrm{Trtria}}
 &\simeq\text{ens des triangles de trièdres parallèles ( ``gds Triangles}\\
 &\simeq\text{ens des pavés \unsure{coorientés|orientés}}
\end{aligned}
\qquad \simeq\text{Ens des \unsure{tribios|tribos|faisceaux réduits}}
\]

% middle left
\diagram{\begin{tikzcd}[ampersand replacement=\&,column sep=large,row sep=large,execute at end picture={\draw ([xshift=-3pt,yshift=2pt]tl.north west) rectangle ([xshift=3pt,yshift=-2pt]bl.south east);}]
|[alias=tl]| \mathrm{tria}'{}^{\unsure{\simeq\mathrm{bibo}}}\arrow[d,"3"'] \& \mathrm{tria}\simeq\mathrm{hex}\arrow[l,"2"']\arrow[d,"3"]\\
\mathrm{Hex}\arrow[d,"3"] \& \deleted{\unsure{\mathrm{Tr}}}\mathrm{Trtria}\deleted{\unsure{\simeq\mathrm{Hex}}}\arrow[l,"2"']\arrow[ul,phantom,"{\text{(cart)}}",very near end]\\
|[alias=bl]| \mathrm{triHex} \&
\end{tikzcd}}

% middle right, beside diagram
\margin{\uline{NB} les rev. d'ordre $2$ sont les passages aux trièdres opposés resp. aux triangles de trièdres par. opposés.}

% lower half, lines 1–3
\uline{NB} Ce diagramme se déduit (à isom can près) du diagramme

% lower half, line 2
\diagram{\begin{tikzcd}[ampersand replacement=\&]
\mathrm{Tria}'\arrow[r]\&\mathrm{Hex}\arrow[r]\&\mathrm{triHex}
\end{tikzcd}}

% lower half, lines 3–7
Car le revêtement $\mathrm{Trtria}$ d'ordre $2$ de $\mathrm{Hex}$ se déduit (mod iso can) du rev. d'ordre $3$ $\mathrm{Hex}$ de $\mathrm{triHex}$, comme le revêtement principal associé. D'autre part, ce diagramme

% lower half, line 8
\diagram{\begin{tikzcd}[ampersand replacement=\&]
\mathrm{Tria}'\arrow[r,"3"]\&\mathrm{Hex}\arrow[r,"3"]\&\mathrm{triHex}
\end{tikzcd}}

% lower half, lines 9–11
permet de récupérer, \deleted{la \unsure{façon|famille|groupe} \unsure{canonique}} \inserted{pour $\tau\in\mathrm{triHex}$,} $\Delta$ via la fibre $(\mathrm{Tria}')_{\tau}^{\unsure{+}}$, considérée comme ens $(3,3,3)$

% archive sheet number: 72
